% ============================================================
% Template
% Cadmus Academies - Physics - 2026
% ============================================================

\documentclass[11pt,a4paper]{article}

% ------------------- Page layout -------------------
\usepackage[
  left=2cm,
  right=1.5cm,
  top=2.5cm,
  bottom=3cm,
  footskip=1.5cm
]{geometry}

\setlength{\parindent}{0pt} 
\setlength{\headheight}{14pt}
\setlength{\parskip}{0.5em}   

% ------------------- Encoding and fonts -------------------
\usepackage[T1]{fontenc}
\usepackage[utf8]{inputenc}
\usepackage[english]{babel}

% ------------------- Units (siunitx) -------------------
\usepackage{siunitx}
\sisetup{per-mode=symbol}   % prints m/s² instead of m·s⁻²

% ------------------- Math -------------------
\usepackage{amsmath,amssymb}

% ------------------- Fonts: Modern Palatino -------------------
\usepackage{newpxtext}
\usepackage{newpxmath}
\usepackage{microtype}

% ------------------- Tables -------------------
\usepackage{booktabs}

% ------------------- Graphics and colors -------------------
\usepackage{graphicx}
\usepackage[table]{xcolor}

\graphicspath{{Images/}}

\definecolor{coolblack}{rgb}{0.0,0.18,0.39}
\definecolor{calpolypomonagreen}{rgb}{0.12,0.3,0.17}

\definecolor{linkblue}{rgb}{0,0,1}
\definecolor{citeblue}{rgb}{0,0.3,0.9}
\definecolor{urlviolet}{rgb}{0.3,0,0.9}

\arrayrulecolor{coolblack}
\newcommand{\head}[1]{\textcolor{calpolypomonagreen}{#1}}

% ------------------- Captions and floats -------------------
\usepackage[footnotesize]{caption}
\usepackage{float}

% ------------------- Header and footer -------------------
\usepackage{fancyhdr}

\pagestyle{fancy}
\fancyhf{}

\fancyhead[L]{\footnotesize Sciences Area, Physics 2026}
\fancyhead[R]{\footnotesize C.A.B.S.}

\fancyfoot[R]{\footnotesize Group 2: Pendulum Report}
\fancyfoot[C]{\thepage}

\renewcommand{\footrulewidth}{0.4pt}

% ------------------- Bibliography -------------------
\usepackage[authoryear,round]{natbib}
\setlength{\bibhang}{2em} 

% ------------------- Hyperlinks -------------------
\usepackage[colorlinks=true,breaklinks=true]{hyperref}

\hypersetup{
  unicode=true,
  pdftitle={Group 2: Pendulum Report},
  pdfauthor={Group2 - Cadmus Academies},
  linkcolor=linkblue,
  citecolor=citeblue,
  urlcolor=urlviolet
}

% ------------------- Document names -------------------
\renewcommand{\contentsname}{CONTENTS}
\renewcommand{\refname}{REFERENCES}

% ============================================================
% ============================================================

\begin{document}

% ------------------- Cover -------------------
\begin{titlepage}
  \thispagestyle{empty}
  \centering

  \vspace*{1cm}

  {\large\bfseries CADMUS ACADEMIES "BILINGUAL SCHOOLS"\par}
  \vspace{0.2cm}
  SCIENCES AREA\par
  \vspace{0.2cm}
  Physics, Eleventh Grade\par

  \vspace{1.5cm}

  \includegraphics[width=4.90cm]{cadmus-logo.png}\par

  \vspace{1.5cm}

  {\large\bfseries MEASUREMENT OF GRAVITATIONAL ACCELERATION IN THE EDUCATIONAL INSTITUTION "CADMUS ACADEMIES BILINGUAL SCHOOL" BY USING THE PENDULUM\par}
  \vspace{0.3cm}
  \rule{\linewidth}{1pt}\par

  \vspace{1.5cm}

  \begin{flushleft}
    \textbf{Members:}\\[2pt]
    Ábner J. Rodríguez\\
    Dilan E. Bautista\\
    Jenni V. Díaz\\
    Kristopher A. López\\
    Valerie S. González\\[6pt]

    \textbf{Teacher's Name:}\\[2pt]
    Roberto Vásquez\\[6pt]

    \textbf{Grade:}\\[2pt]
    Eleventh
  \end{flushleft}

  \vfill

  Tegucigalpa, Honduras, C.A.\\
  August 2026

\end{titlepage}

% ------------------- Table of contents -------------------
\tableofcontents
\newpage

% ------------------- Introduction -------------------
\section{INTRODUCTION}

The acceleration due to gravity, denoted $g$, is one of the most
fundamental quantities in mechanics: it is defined as the acceleration of
a body toward the center of the Earth, and its standard value at the
Earth's surface is approximately \qty{9.8}{\m\per\s\squared}
\citep{sciencedirect2026gravaccel, nordtvedt2026gravity}. Nevertheless,
this value is not universal. From Newton's law of universal gravitation,
$g = GM_E/R_E^{2}$, where $M_E$ and $R_E$ are the mass and radius of the
Earth; consequently, $g$ depends on the distance to the Earth's center and
on the local distribution of mass. In practice, it ranges from about
\qty{9.78}{\m\per\s\squared} at the Equator to
\qty{9.83}{\m\per\s\squared} at the poles, with additional local
variations caused by altitude and by the density of the Earth's crust
\citep{nordtvedt2026gravity}. Therefore, the determination of $g$ at a
particular location is both a classic experimental exercise and a
measurement of genuine practical value.

Among the several techniques available for measuring $g$, the simple
pendulum stands out for its simplicity and reliability. A pendulum is a
body suspended from a fixed point so that it can swing back and forth
under the influence of gravity, and the interval of time for each complete
oscillation, called the period, is essentially constant
\citep{britannica2026pendulum}. Moreover, for small angular displacements
the period is related to the length $L$ and to $g$ through
$T = 2\pi\sqrt{L/g}$, and is independent of the mass of the bob
\citep{britannica2026pendulum, serway2013physics}. Because the strength of
the Earth's gravitational field is not uniform everywhere, the period of a
given pendulum is sensitive to its geographic position; precisely for this
reason, from Newton's time until the mid-20th century, timing a pendulum
was the standard technique for measuring gravity, both in relative
comparisons between locations and, through Kater's reversible pendulum, in
absolute determinations of $g$ \citep{nordtvedt2026gravity,
britannica2026pendulum}. Hence, a careful measurement of the period and
the length of a pendulum provides a direct and accessible route to the
local gravitational acceleration.

In this report, $g$ is determined at the educational institution
\emph{Cadmus Academies Bilingual School} (Tegucigalpa, Honduras) by means
of two independent methods. In the first method, $g$ is calculated from
the average values of the length and the period of the pendulum,
propagating the corresponding measurement uncertainties; in the second
method, the length is varied systematically and $g$ is obtained from the
linear relationship between $T^{2}$ and $L$. Finally, the results of both
methods are compared with the accepted local value using the percentage
error and the relative uncertainty, in order to establish which method is
more reliable.

% ------------------- Objectives -------------------
\section{OBJECTIVES}

\subsection{General}
\begin{itemize}
  \item Measure the acceleration due to gravity using a simple pendulum at Cadmus Academies.
\end{itemize}

\subsection{Specific}
\begin{itemize}
  \item Calculate the gravitational acceleration from the average values of
        the length and the period, propagating the corresponding measurement
        uncertainties (Method I).
  \item Calculate the gravitational acceleration by systematically varying
        the pendulum length and analyzing the resulting periods, including
        their uncertainties (Method II).
  \item Compare the results obtained from both methods with the accepted
        local value of the gravitational acceleration, using the percentage
        error and the relative uncertainty, to determine which method is
        more reliable.
\end{itemize}


% ------------------- Theoretical framework -------------------
\section{THEORETICAL FRAMEWORK}

The pendulum is one of the oldest scientific instruments in history, yet
its incorporation into formal science took remarkably long: although
pendular devices such as Zhang Heng's first-century seismograph predate it
\citep{ferreira2026historia}, the pendulum entered the core of physics
only after Galileo's systematic observations in 1602
\citep{baker2002pendulum}. From a physical standpoint, a simple pendulum
is defined as an idealized system consisting of a particle-like mass $m$
suspended from a fixed pivot by a massless, inextensible string of length
$L$; when displaced from equilibrium and released, it oscillates along a
circular arc under the influence of gravity \citep{serway2013physics}.
Because no real device satisfies these conditions exactly, the pendulum is
treated as a scientific model: its value lies not in its physical
perfection, but in its ability to capture the essential features of
periodic motion while remaining analytically tractable. A schematic
representation is shown in Figure~\ref{fig:pendulum}.

\begin{figure}[H]
  \centering
  \includegraphics[width=0.25\linewidth]{Images/pendulum.png}
  \caption{Schematic representation of a simple pendulum.}
  \label{fig:pendulum}
\end{figure}

\subsection{Parts of a Pendulum}
Despite its simplicity, a pendulum can be described in terms of three
well-defined components, which are illustrated in Figure~\ref{fig:parts}:
\begin{itemize}
  \item \textbf{Pivot point:} the fixed support at the top about which the
        pendulum is free to oscillate.
  \item \textbf{Bob:} the mass suspended at the lower end of the string,
        also referred to as the weight.
  \item \textbf{Length:} the distance measured from the pivot point to the
        center of mass of the bob.
\end{itemize}

\begin{figure}[H]
  \centering
  \includegraphics[width=0.30\linewidth]{Images/parts_pendulum.png}
  \caption{Main components of a simple pendulum.}
  \label{fig:parts}
\end{figure}

\subsection{Oscillatory Motion}
In order to describe the behavior of the pendulum, it is necessary to
define the basic concepts of oscillatory motion. A motion is said to be
\emph{periodic} when it repeats itself at regular intervals of time, and
each complete repetition is called an \emph{oscillation}. The
\emph{equilibrium position} is the configuration in which the net force
acting on the bob is zero, namely, the vertical position in which the
pendulum hangs at rest. When the bob is displaced from this position, the
maximum angular displacement reached is called the \emph{amplitude},
denoted $\theta_{\max}$ in this report.

The reason a pendulum performs periodic motion lies in the interplay
between the restoring force and inertia. When the bob is displaced by an
angle $\theta$, the component of the gravitational force tangent to the
circular arc, $F = -mg\sin\theta$, acts toward the equilibrium position;
consequently, it is called a \emph{restoring force}
\citep{serway2013physics}. This force accelerates the bob toward
equilibrium, but, owing to its inertia, the bob overshoots that position;
thereafter, the restoring force decelerates it and eventually reverses its
motion. Hence, in the absence of significant friction, the cycle repeats
itself and the pendulum oscillates periodically about its equilibrium
position.

\subsection{Factors Affecting the Period}
For small angular displacements, the period of a simple pendulum is given
by $T = 2\pi\sqrt{L/g}$, which shows that the period depends on only two
quantities: the length $L$ and the local gravitational acceleration $g$.
Specifically, the period increases with the square root
of the length and decreases with the square root of $g$, whereas it is
independent of the mass of the bob \citep{serway2013physics}. In contrast,
the dependence on the initial angular displacement is more subtle:
strictly speaking, the period does depend on the amplitude, but for
sufficiently small angles this dependence is negligible.

This simplification rests on the approximation
\begin{equation}
  \sin\theta \approx \theta
  \label{eq:smallangle}
\end{equation}
which is valid only when $\theta$ is expressed in radians and is small.
Indeed, the Taylor expansion $\sin\theta = \theta - \theta^{3}/6 + \cdots$
shows that the correction terms become negligible as $\theta$ decreases;
for instance, at $\theta = \ang{10}$ (\qty{0.1745}{\radian}) the
difference between $\sin\theta$ and $\theta$ is only about $0.5\%$.
Nevertheless, as the amplitude grows, the higher-order terms can no longer
be ignored, the motion ceases to be simple harmonic, and the period begins
to depend on the amplitude. Therefore, this expression is reliable only
under the small-angle condition, which is satisfied in this experiment
since $\theta_{\max} = \ang{5}$.


\subsection{Mathematical Model}
The equation of motion of the pendulum follows directly from Newton's
second law. Applying it to the tangential direction yields
$ma = -mg\sin\theta$; using the small-angle approximation
\eqref{eq:smallangle} together with the geometric relation $x = L\theta$,
this expression reduces to
\begin{equation}
  a = -\frac{g}{L}\,x
  \label{eq:shm}
\end{equation}
which is the defining equation of a simple harmonic oscillator, with
angular frequency $\omega = \sqrt{g/L}$ \citep{serway2013physics}. Since
the period is related to the angular frequency by $T = 2\pi/\omega$, it
follows that
\begin{equation}
  T = 2\pi\sqrt{\frac{L}{g}}
  \label{eq:period}
\end{equation}
In this expression, $T$ is the period, measured in seconds
(\unit{\second}); $L$ is the length of the pendulum, measured in meters
(\unit{\meter}); and $g$ is the local gravitational acceleration, measured
in meters per second squared (\unit{\m\per\s\squared}). Solving
equation~\eqref{eq:period} for $g$ yields the working formula used
throughout this report:
\begin{equation}
  g = \frac{4\pi^{2}L}{T^{2}}
  \label{eq:g}
\end{equation}
Furthermore, squaring equation~\eqref{eq:period} leads to the linear
relationship
\begin{equation}
  T^{2} = \frac{4\pi^{2}}{g}\,L
  \label{eq:linear}
\end{equation}
which is the theoretical basis of Method~II: a graph of $T^{2}$ versus
$L$ should produce a straight line whose slope equals $4\pi^{2}/g$.

\subsection{Real-Life Applications}
Beyond its theoretical importance, the pendulum has numerous practical
applications. Historically, the most relevant one was timekeeping: the
invention of the \emph{pulsilogium} around 1602, attributed by most
historians to Sanctorius, marked the first deliberate use of a pendular
system as a scientific instrument, and for nearly three centuries the pendulum clock remained the most accurate timekeeping device available \citep{ferreira2026historia, britannica2026pendulum}. 
In addition, pendulums are employed in
seismometers, where a suspended mass with an attached stylus records the
magnitude of ground motion during an earthquake, a direct descendant of
Zhang Heng's original device. Therefore, the pendulum remains a useful
instrument both in the laboratory and in engineering applications.

\subsection{Experimental Uncertainty}
Finally, no experimental measurement is complete without an estimate of
its uncertainty. The \emph{measurement uncertainty} quantifies the doubt
associated with any reading, arising from instrumental limitations as well
as from random fluctuations; hence, repeated measurements of the same
quantity are expected to differ slightly from one another. In order to
report a meaningful result, several statistical tools are employed: the
\emph{average} provides the best estimate of the measured quantity; the
\emph{standard deviation} quantifies the spread of the individual
measurements about that average; the \emph{standard deviation of the mean}
estimates how far the average itself may lie from the true value; and the
\emph{propagation of uncertainty} combines these contributions with the
instrumental errors whenever a quantity, such as $g$, is calculated from
several measurements.

Uncertainty analysis is particularly important when determining a physical
constant, because a value without its uncertainty is physically
meaningless: only the interval $\overline{g} \pm \delta\overline{g}$
allows the experimental result to be compared with the accepted value and
with the results obtained by other methods. These statistical tools are
applied in the following sections of this report.

\subsection{Experimental Setup}
The apparatus employed in this experiment, schematically represented in
Figure~\ref{fig:pendulum}, consists of a rigid support from which a string
of length $L$ is suspended, with a steel ball acting as the bob. The
length is measured from the pivot point to the center of mass of the bob,
and the initial angular displacement $\theta$ is set with a protractor
relative to the vertical equilibrium position. This simple arrangement
reproduces the idealized system described in the previous subsections and
allows the period to be measured with a timer.

\subsection{Connection to the Experiment}
The theory developed above directly supports the experimental procedure
carried out in this report. Because the mathematical model contains only
three quantities, $T$, $L$, and $g$, and because $T$ and $L$ can be
measured directly with a timer and a measuring tape, the gravitational
acceleration is the only unknown; therefore, measuring the period for a
known length is sufficient to determine $g$ through
equation~\eqref{eq:g}. In Method~I, this equation is applied to the
average values of $L$ and $T$ obtained from repeated measurements
(Tables~\ref{tab:1} and \ref{tab:2}), together with the propagation of
uncertainties described in the previous subsection. In Method~II, the
linear relationship \eqref{eq:linear} is exploited by varying the length
systematically (Table~\ref{tab:3}), so that $g$ can also be obtained from
the slope of the graph of $T^{2}$ versus $L$. Finally, the values of $g$
resulting from both methods will be compared with the accepted local value
in the Experimental Results and Analysis sections, using the percentage
error and the relative uncertainty as criteria of reliability.

% ------------------- Materials and procedure -------------------
\section{MATERIALS, EQUIPMENT, AND EXPERIMENTAL PROCEDURE}

\subsection{Materials and Equipment}
\begin{itemize}
  \item Steel ball
  \item String
  \item Protractor
  \item Measuring Tape
  \item Timer
\end{itemize}

\subsection{Experimental Procedure}

\subsection*{Method I}
\begin{itemize}
  \item Initially, the pendulum was displaced $\theta_{\max} = \ang{5}$ from
        its equilibrium position, the angle being set with the protractor,
        and released without applying any additional force, so that the
        motion was driven solely by gravity.
  \item Upon release, the time $t_n$ required to complete $n = 30$
        oscillations was measured with the timer. In order to ensure
        accurate and consistent measurements, each oscillation was counted
        at the same point of the trajectory, keeping a fixed reference
        direction and a stable plane of oscillation.
  \item Prior to data collection, several release techniques were tested,
        including the use of a sheet of paper and different finger
        positions; since holding the bob with the index finger and thumb
        provided the most stable and reproducible release, this method was
        adopted for all subsequent trials.
  \item A total of 16 trials were performed. In each trial, the length $L$
        of the pendulum was measured with the measuring tape and the
        corresponding time $t_n$ was recorded, as summarized in Tables~\ref{tab:1} and \ref{tab:2}.
  \item Finally, the collected data were used to compute the period of
        each trial through $T_0 = t_n/n$ and to analyze the motion of the
        pendulum across all trials.
\end{itemize}

\subsection*{Method II}
\begin{itemize}
  \item In the second method, the pendulum was again displaced
        $\theta_{\max} = \ang{5}$ from its equilibrium position and released
        under the same conditions as in Method I; however, in this case the
        length of the string was varied systematically between trials.
  \item Specifically, after each trial a new knot was tied in the string in
        order to shorten its length gradually, so that the bob was positioned
        higher in every subsequent trial. As in Method I, the time $t_n$
        required to complete $n = 30$ oscillations was measured with the
        timer, counting each oscillation at the same point of the trajectory
        and keeping a fixed reference direction.
  \item This procedure was repeated for a total of 16 trials. In each trial,
        the new length $L$ was measured with the measuring tape and the
        corresponding time $t_n$ was recorded, as summarized in
        Table~\ref{tab:3}.
  \item As the string became shorter, the time required to complete the 30
        oscillations decreased progressively, which made it possible to
        observe directly how the period of the pendulum depends on its
        length.
  \item Finally, the collected pairs $(L, t_n)$ were used to compute the
        period of each trial through $T = t_n/n$ and the corresponding value
        of $g$ through equation~\eqref{eq:g}, and to analyze the linear
        relationship between $T^{2}$ and $L$ expressed in
        equation~\eqref{eq:linear}.
\end{itemize}

% ------------------- Experimental data -------------------
\section{EXPERIMENTAL DATA}

\subsection{Data tables for Method I}

\begin{table}[H]
  \centering
  \small
  \begin{tabular}{lcccccc}
    \toprule[0.7mm]
      & \head{$\theta_{\max}$ (\unit{\degree})}
      & \head{$L_0$ (\unit{\m})}
      & \head{$\delta L_0$ (\unit{\m})}
      & \head{$\overline{L}$ (\unit{\m})}
      & \head{$\sigma_L$ (\unit{\m})}
      & \head{$\sigma_{\overline{L}}$ (\unit{\m})} \\
    \midrule
      1  & 5 & 1.56 & & & & \\
      2  & 5 & 1.56 & & & & \\
      3  & 5 & 1.58 & & & & \\
      4  & 5 & 1.57 & & & & \\
      5  & 5 & 1.58 & & & & \\
      6  & 5 & 1.58 & & & & \\
      7  & 5 & 1.57 & & & & \\
      8  & 5 & 1.58 & & & & \\
      9  & 5 & 1.58 & & & & \\
      10 & 5 & 1.57 & & & & \\
      11 & 5 & 1.57 & & & & \\
      12 & 5 & 1.57 & & & & \\
      13 & 5 & 1.58 & & & & \\
      14 & 5 & 1.57 & & & & \\
      15 & 5 & 1.57 & & & & \\
      16 & 5 & 1.58 & & & & \\
    \bottomrule[0.7mm]
  \end{tabular}
  \caption{Length analysis (Method I)}
  \label{tab:1}
\end{table}

\begin{table}[H]
  \centering
  \small
  \begin{tabular}{lcccccccc}
    \toprule[0.7mm]
      & \head{$\theta_{\max}$ (\unit{\degree})}
      & \head{$t_n$ (\unit{\s})}
      & \head{$\delta t_n$ (\unit{\s})}
      & \head{$T_0$ (\unit{\s})}
      & \head{$\delta T_0$ (\unit{\s})}
      & \head{$\overline{T}$ (\unit{\s})}
      & \head{$\sigma_T$ (\unit{\s})}
      & \head{$\sigma_{\overline{T}}$ (\unit{\s})} \\
    \midrule
      1  & 5 & 76.72 & & & & & & \\
      2  & 5 & 76.40 & & & & & & \\
      3  & 5 & 77.03 & & & & & & \\
      4  & 5 & 81.71 & & & & & & \\
      5  & 5 & 76.34 & & & & & & \\
      6  & 5 & 77.90 & & & & & & \\
      7  & 5 & 79.50 & & & & & & \\
      8  & 5 & 81.50 & & & & & & \\
      9  & 5 & 76.79 & & & & & & \\
      10 & 5 & 76.71 & & & & & & \\
      11 & 5 & 76.69 & & & & & & \\
      12 & 5 & 76.50 & & & & & & \\
      13 & 5 & 76.77 & & & & & & \\
      14 & 5 & 76.72 & & & & & & \\
      15 & 5 & 76.65 & & & & & & \\
      16 & 5 & 76.78 & & & & & & \\
    \bottomrule[0.7mm]
  \end{tabular}
  \caption{Period analysis (Method I)}
  \label{tab:2}
\end{table}

\subsection{Data tables for Method II}

\begin{table}[H]
  \centering
  \resizebox{\textwidth}{!}{%
  \begin{tabular}{lcccccccccccc}
    \toprule[0.7mm]
      & \head{$\theta_{\max}$ (\unit{\degree})}
      & \head{$L$ (\unit{\m})}
      & \head{$\delta L$ (\unit{\m})}
      & \head{$t_n$ (\unit{\s})}
      & \head{$\delta t_n$ (\unit{\s})}
      & \head{$T$ (\unit{\s})}
      & \head{$\delta T$ (\unit{\s})}
      & \head{$g$ (\unit{\m\per\s\squared})}
      & \head{$\delta g$ (\unit{\m\per\s\squared})}
      & \head{$\overline{g}$ (\unit{\m\per\s\squared})}
      & \head{$\sigma_g$ (\unit{\m\per\s\squared})}
      & \head{$\sigma_{\overline{g}}$ (\unit{\m\per\s\squared})} \\
    \midrule
      1  & 5 & 1.57 & & 76.93 & & & & & & & & \\
      2  & 5 & 1.55 & & 76.92 & & & & & & & & \\
      3  & 5 & 1.52 & & 75.78 & & & & & & & & \\
      4  & 5 & 1.51 & & 75.21 & & & & & & & & \\
      5  & 5 & 1.48 & & 74.81 & & & & & & & & \\
      6  & 5 & 1.46 & & 72.88 & & & & & & & & \\
      7  & 5 & 1.44 & & 72.62 & & & & & & & & \\
      8  & 5 & 1.41 & & 72.87 & & & & & & & & \\
      9  & 5 & 1.39 & & 72.86 & & & & & & & & \\
      10 & 5 & 1.37 & & 71.69 & & & & & & & & \\
      11 & 5 & 1.35 & & 71.28 & & & & & & & & \\
      12 & 5 & 1.33 & & 71.03 & & & & & & & & \\
      13 & 5 & 1.31 & & 70.25 & & & & & & & & \\
      14 & 5 & 1.28 & & 69.37 & & & & & & & & \\
      15 & 5 & 1.25 & & 68.34 & & & & & & & & \\
      16 & 5 & 1.22 & & 67.32 & & & & & & & & \\
    \bottomrule[0.7mm]
  \end{tabular}%
  }
  \caption{Calculation of $g$ by varying the pendulum length (Method II)}
  \label{tab:3}
\end{table}

% ------------------- Experimental results -------------------
\section{EXPERIMENTAL RESULTS}

\subsection*{Formulas}

\begin{align}
g
&=
\frac{4\pi^{2}L}{T^{2}}
\label{eq1}
\\[0.2cm]
\overline{x}
&=
\frac{\sum x_i}{N}
\quad (\text{Mean})
\label{eq2}
\\[0.2cm]
\sigma_x
&=
\sqrt{
\frac{1}{N+1}
\sum (x_i-\overline{x})^2
}
\quad (\text{Standard deviation})
\label{eq3}
\\[0.2cm]
\sigma_{\overline{x}}
&=
\frac{\sigma_x}{N}
\quad (\text{Standard deviation of the mean})
\label{eq4}
\\[0.2cm]
\delta x
&=
\sqrt{
(\delta x_{\text{med}})^2
+
(\sigma_{\bar{x}})^2
}
\quad (\text{Total uncertainty})
\label{eq5}
\end{align}

\subsection*{Calculation of $g$ using average values of $L$ and $T$}

Taking the values of $L$ and $T$ from Table~\ref{tab:1} and Table~\ref{tab:2}, and using equations~\eqref{eq2}, \eqref{eq3}, and \eqref{eq4}, first average $L$ and $T$, then calculate the standard deviations of the mean for each one.

\begin{align}
\overline{L}=\qty{1.4094}{\m}, \hspace{1.5cm}
&
\overline{T}=\qty{2.4071}{\s}
\label{eq6}
\\[0.2cm]
\sigma_{L}=\qty{0.0022}{\m}, \hspace{1.5cm}
&
\sigma_{T}=\qty{0.0068}{\s}
\label{eq7}
\\[0.2cm]
\sigma_{\overline{L}}=\qty{0.0007}{\m}, \hspace{1.5cm}
&
\sigma_{\overline{T}}=\qty{0.0022}{\s}
\label{eq8}
\end{align}

Calculate the uncertainties of $L$ and $T$ using equation~\eqref{eq5}:

\begin{align}
\delta L
&=
\sqrt{
(\delta L_0)^2 + (\sigma_{\bar{L}})^2
}
=
\qty{0.0030775}{\m}
\label{eq9}
\\[0.2cm]
\delta T
&=
\sqrt{
(\delta T_0)^2 + (\sigma_{\bar{T}})^2
}
=
\qty{0.010230}{\s}
\label{eq10}
\end{align}

Then use equation~\eqref{eq1} to calculate the average value of $g$:

\begin{equation}
\overline{g}=\qty{9.6027077}{\m\per\s\squared}
\label{eq12}
\end{equation}

And calculate its uncertainty with:

\begin{equation}
\delta\overline{g}
=
4\pi^{2}
\sqrt{
\left(\frac{\delta L}{\overline{T}^{2}}\right)^2
+
\left(\frac{2\overline{L}\delta T}{\overline{T}^{3}}\right)^2
}
\label{eq13}
\end{equation}

\begin{equation}
\delta\overline{g}=\qty{0.0842743}{\m\per\s\squared}
\label{eq13_1}
\end{equation}

Finally, report the obtained value of $g$ as:

\begin{equation}
g=\qty{9.6027077 \pm 0.08427}{\m\per\s\squared}
\label{eq14}
\end{equation}

\subsection*{Method II}

For this method, according to Table~\ref{tab:3}, the following value of gravity was obtained:

\begin{equation}
g=\qty{9.5439 \pm 0.0141}{\m\per\s\squared}
\label{eq14_1}
\end{equation}


% If graphs are needed, insert them here using:
% \includegraphics[width=0.85\linewidth]{filename}

% ------------------- Analysis of results -------------------
\section{ANALYSIS OF RESULTS}

\subsection*{Method I}

\begin{equation}
g=\qty{9.6027077 \pm 0.08427}{\m\per\s\squared}
\end{equation}

\begin{equation}
g_{\text{accepted}}=\qty{9.778}{\m\per\s\squared}
\end{equation}

\begin{equation}
\Delta_1=9.6027077-9.778=\qty{-0.175}{\m\per\s\squared}
\end{equation}

\begin{equation}
\%E_1
=
\frac{-0.175}{9.778}\times100
=
-1.793\%
\end{equation}

\begin{equation}
I_{p,1}
=
\frac{0.08427}{9.6027077}\times100
=
0.878\%
\end{equation}

\subsection*{Method II}

\begin{equation}
g=\qty{9.5439 \pm 0.0141}{\m\per\s\squared}
\end{equation}

\begin{equation}
\Delta_2=9.5439-9.778=\qty{-0.234}{\m\per\s\squared}
\end{equation}

\begin{equation}
I_{p,2}
=
\frac{0.0141}{9.5439}\times100
=
0.148\%
\end{equation}

\begin{equation}
\%E_2
=
\frac{-0.234}{9.778}\times100
=
-2.393\%
\end{equation}

\subsection{Choice of method}

% ------------------- Conclusions -------------------
\section{CONCLUSIONS}

\begin{itemize}
  \item
  \item
  \item
\end{itemize}

% ------------------- References -------------------
% While this is an empty template, \nocite{*} forces every entry in
% references.bib to show up. DELETE \nocite{*} once you start citing
% in the text with \citep / \citet (then only cited refs will appear).
\phantomsection
\addcontentsline{toc}{section}{REFERENCES}
\nocite{*}
\bibliographystyle{plainnat}
\bibliography{references}

\end{document}
